%!TEX root = ../../main.tex
\begin{table}[htbp]
\centering
\small
\setlength{\tabcolsep}{5pt}
\begin{threeparttable}
\caption{Post $\times$ high exposure and the new-hire share, all cells}
\label{tab:newhire}
\begin{tabular}{lcrrr}
\toprule
 & Coefficient (SE) & New-hire share & Cells & Clusters \\
 &  & 2022, \% &  &  \\
\midrule
\addlinespace
\multicolumn{5}{@{}l}{\textit{Pooled, 7 countries with a tenure question}} \\
Post $\times$ High exposure ($\delta$) & $-0.0072^{***}$ $(0.0024)$ & 21.9 & 40\,154 & 485 \\
\addlinespace
\multicolumn{5}{@{}l}{\textit{By country}} \\
Argentina & $0.0018$ $(0.0046)$ & 15.9 & 2\,835 & 31 \\
Brazil & $-0.0111^{***}$ $(0.0032)$ & 24.0 & 14\,960 & 111 \\
Colombia & $-0.0074$ $(0.0062)$ & 25.9 & 8\,892 & 101 \\
Ecuador & $0.0075^{**}$ $(0.0037)$ & 8.3 & 5\,415 & 82 \\
Mexico & $0.0024$ $(0.0071)$ & 17.9 & 1\,932 & 39 \\
South Africa & $0.0086$ $(0.0081)$ & 17.7 & 1\,836 & 46 \\
Uruguay & $-0.0131^{*}$ $(0.0074)$ & 13.9 & 4\,284 & 75 \\
\bottomrule
\end{tabular}
\begin{tablenotes}[flushleft]
\footnotesize
\item \textit{Notes:} Difference in differences on the cell's new-hire share, the weighted fraction of its employed persons who started their current job within the previous twelve months, in share points (multiply by 100 for percentage points): $s_{it} = \delta\,\text{Post}_t \times \text{High}_{o} + \alpha_i + \gamma_{c,a,s,t} + \varepsilon_{it}$, where High marks occupations in the top employment-weighted quintile of the ILO generative-AI exposure score at the country's ISCO digit level. Post $=1$ from 2023Q1 onward (reference 2022Q4). Cell fixed effects and country $\times$ age group $\times$ sex $\times$ quarter fixed effects; weights are mean 2022 cell employment. Sample: Argentina, Brazil, Colombia, Ecuador, Mexico, South Africa, Uruguay, the countries whose survey asks job tenure; Uruguay asks it at first interviews only and Mexico in first-interview quarters only, so their shares rest on part of each cell. Country rows re-estimate the model within the country. New-hire share 2022 is the employment-weighted mean of the outcome over the sample's 2022 cell-quarters. Standard errors in parentheses, clustered by country $\times$ occupation; $^{*}$ $p<0.10$, $^{**}$ $p<0.05$, $^{***}$ $p<0.01$ from the cluster-robust $t$ statistic with $G-1$ degrees of freedom.
\end{tablenotes}
\end{threeparttable}
\end{table}
